\documentclass[12pt]{article}
\usepackage[T1]{fontenc}
\usepackage[utf8]{inputenc}
\usepackage{lmodern}
\usepackage{textcomp}
\usepackage{enumitem}
\usepackage{geometry}
\geometry{margin=1in}
\begin{document}
\begin{center}
Answer Key - Chemistry - Undergraduate Level Assessment 22 \
\small (Answers are ordered exactly as in the question paper.)
\end{center}

\section*{Multiple Choice}
\begin{enumerate}[label=\arabic*.]
\item \textbf{Answer:} B
\\ \textit{Options:} A) It is a colorless gas.; B) It is a red-brown volatile liquid.; C) It is a colorless volatile liquid.; D) It is a yellow metallic solid.
\\ \textit{Note:} At standard temperature and pressure, bromine (Br$_{2}$) exists as a red-brown volatile liquid due to its molecular structure and properties, distinguishing it from other halogens which may exist in different states under the same conditions.
\item \textbf{Answer:} A
\\ \textit{Options:} A) It has little effect.; B) It significantly increases the pH.; C) It significantly decreases the pH.; D) It changes the pH asymptotically to the pKa of the acid.
\\ \textit{Note:} Doubling the volume of the buffer solution by adding water dilutes both the weak acid and its conjugate base equally, thereby maintaining the ratio of their concentrations and resulting in little effect on the pH.
\item \textbf{Answer:} C
\\ \textit{Options:} A) 101.1 x $10^7$ T-1 s-1; B) 26.75 x $10^7$ T-1 s-1; C) 7.081 x $10^7$ T-1 s-1; D) 0.9377 x $10^7$ T-1 s-1
\\ \textit{Note:} The magnetogyric ratio is calculated by dividing the NMR frequency of the nucleus by the magnetic field strength of the NMR spectrometer, and in this case, 198.52 MHz divided by 750 MHz, when converted to the appropriate units, yields a magnetogyric ratio of 7.081 x $10^7$ T-1 s-1.
\item \textbf{Answer:} A
\\ \textit{Options:} A) carbonic acid; B) a halogenated hydrocarbon; C) lead; D) mercury
\\ \textit{Note:} Carbonic acid, which forms when carbon dioxide dissolves in water, is a naturally occurring substance in the carbon cycle and is not classified as a toxic pollutant because it is present in balanced amounts in aquatic ecosystems and plays a role in regulating pH levels.
\item \textbf{Answer:} D
\\ \textit{Options:} A) Ammonium hydroxide; B) Sulfuric acid; C) Acetic acid; D) Potassium hydrogen phthalate
\\ \textit{Note:} Potassium hydrogen phthalate is a primary standard for standardizing bases because it is a stable, non-hygroscopic, and highly pure acid that reacts with bases in a predictable and quantifiable manner, allowing for accurate titration results.
\end{enumerate}

\section*{True / False}
\begin{enumerate}[label=\arabic*.]
\setcounter{enumi}{5}
\item \textbf{Answer:} True
\\ \textit{Options:} A) True; B) False
\\ \textit{Note:} H$_{3}$O+ is formed when H$_{2}$O donates a proton (H+), making H$_{3}$O+ the conjugate acid and H$_{2}$O the corresponding conjugate base in this acid-base pair.
\item \textbf{Answer:} True
\\ \textit{Options:} A) True; B) False
\\ \textit{Note:} The molecular geometry of thionyl chloride (SOCl 2) is trigonal pyramidal because it has one lone pair of electrons on the sulfur atom, which affects the spatial arrangement of the three bonding pairs of chlorine and oxygen, resulting in a geometry that is not planar.
\item \textbf{Answer:} False
\\ \textit{Options:} A) True; B) False
\\ \textit{Note:} The infrared-active normal modes of a carbon dioxide molecule include asymmetric stretching and bending, but not symmetric stretching, as the latter does not result in a change in the dipole moment necessary for infrared activity.
\item \textbf{Answer:} False
\\ \textit{Options:} A) True; B) False
\\ \textit{Note:} The solid-state structures of the principal allotropes of elemental boron are not composed of B$_{8}$ cubes, as they primarily consist of complex icosahedral clusters and various interconnected polyhedral structures instead.
\item \textbf{Answer:} True
\\ \textit{Options:} A) True; B) False
\\ \textit{Note:} The amide ion (NH 2-) is considered the strongest base in liquid ammonia because it has a high affinity for protons, making it more capable of accepting protons than any other species present in that solvent.
\end{enumerate}

\section*{Long Form Response}
\begin{enumerate}[label=\arabic*.]
\setcounter{enumi}{10}
\item \textbf{Answer:} The relationship between pH and chemical shifts in the ^31 P NMR spectrum of sodium phosphate is primarily influenced by the protonation state of the phosphate species. When the pH equals the pKa of H$_{2}$PO$_{4}$^-, which is approximately 7.2, there is a significant balance between H$_{2}$PO$_{4}$^- and HPO $4^2$- species. At this point, a notable chemical shift occurs, leading to a resonance signal typically observed at around 4.62 ppm. This shift can be attributed to the differing electronic environments of the phosphorus atom in these two protonation states, which affects its shielding and consequently the chemical shift in the NMR spectrum. Understanding this relationship is crucial for interpreting phosphate chemistry in biological and chemical systems.
\\ \textit{Note:} correct because at the pH equal to the pKa of H$_{2}$PO$_{4}$^-, the equilibrium between H$_{2}$PO$_{4}$^- and HPO $4^2$- alters the electronic environment around the phosphorus atom, resulting in a distinct chemical shift in the ^31 P NMR spectrum, which is observed at around 4.62 ppm.
\item \textbf{Answer:} The pH of 0.1 M aqueous solutions of various salts is influenced by the nature of the ions that result from their dissociation in water. Salts composed of strong acids and strong bases, such as sodium chloride (NaCl), do not hydrolyze and thus do not affect the pH of the solution, resulting in a neutral pH of around 7. In contrast, salts derived from weak acids or bases can lead to hydrolysis, which alters the pH. For example, ammonium chloride (NH 4 Cl) produces an acidic solution, while sodium acetate (NaC 2 H$_{3}$O$_{2}$) creates a basic solution. Therefore, NaCl yields the lowest pH among various salts because it remains neutral, without contributing to any acidic or basic properties in the solution compared to those that do undergo hydrolysis.
\\ \textit{Note:} correct because it accurately identifies that the pH of salt solutions is determined by the hydrolysis of the resulting ions, with NaCl remaining neutral due to its origin from a strong acid and strong base, while salts from weak acids or bases can lead to acidic or basic solutions.
\end{enumerate}

\vfill
\noindent\textit{End of Answer Key}
\end{document}